% ============================================================
%  Chương VII – Thách thức và xu hướng
%  ~400 từ · 1.5 trang
% ============================================================
\chapter{Thách thức và xu hướng}
\label{chap:challenges}

% Ý chính: PKI hiện tại không hoàn hảo — vừa có vấn đề kỹ thuật,
% vừa có vấn đề quản trị (governance). Xu hướng tương lai đang giải quyết từng lớp.

% ----------------------------------------------------------
\section{Thách thức hiện tại}
\label{sec:7.1}

\subsubsection*{CA Compromise — bài học từ DigiNotar và Symantec}
% Ý chính: DigiNotar 2011 \cite{Fox-IT2012}: 531 cert giả, Iran MITM Gmail 300K user
% → DigiNotar phá sản trong vài tuần sau khi bị distrust
% Symantec 2017–2018 \cite{Sleevi2017}: mis-issue hàng nghìn cert; Google/Mozilla
% thực hiện distrust kéo dài 1 năm để không break internet
% Bài học: CA tập trung là single point of failure về trust; cần accountability

\subsubsection*{Certificate Transparency (CT) như giải pháp giám sát}
% Ý chính: append-only Merkle tree log công khai; mọi cert TLS phải được log
% trước khi Chrome/Safari chấp nhận (từ 2018) \cite{RFC9162}
% Domain owner dùng crt.sh, Google CT để monitor cert giả cho domain của mình
% CT không ngăn mis-issuance nhưng đảm bảo mis-issuance bị phát hiện

\subsubsection*{Quản lý vòng đời chứng chỉ}
% Ý chính: cert expiry là rủi ro vận hành thực tế — nhiều outage lớn do quên renew
% CAA DNS records \cite{RFC8659}: whitelist CA được phép issue cho domain
% Thách thức ở enterprise: hàng nghìn cert, nhiều team, thiếu automation

\subsubsection*{PKI cho IoT}
% Ý chính: billions of devices, tài nguyên giới hạn (RAM < 64KB, không có UI)
% Cần: lightweight cert format, device identity provisioning tự động,
% Thách thức: làm sao revoke cert của device bị mất/bị hack?

% ----------------------------------------------------------
\section{Xu hướng mới}
\label{sec:7.2}

\subsubsection*{Post-Quantum PKI}
% Ý chính: Shor's algorithm sẽ phá RSA và ECC trên quantum computer đủ mạnh
% NIST PQC 2024 \cite{NIST-PQC2024}: FIPS 203 (Kyber/CRYSTALS-KEM),
% FIPS 204 (Dilithium), FIPS 205 (SPHINCS+)
% Thách thức migration: cert size lớn hơn nhiều (Dilithium sig ~2.4KB vs ECDSA 64B),
% cần "hybrid" cert (cả ECC + PQC) trong giai đoạn chuyển tiếp

\subsubsection*{Tự động hoá và rút ngắn thời hạn chứng chỉ}
% Ý chính: ACME \cite{RFC8555} + EST → cert lifecycle hoàn toàn tự động
% Apple mandate 2023: max 398 ngày; xu hướng 90 ngày; đích đến: 24 giờ
% Short-lived cert → revocation trở nên ít quan trọng hơn (cert tự hết hạn)

\subsubsection*{Decentralized Identity (DID)}
% Ý chính: W3C DID v1.0 \cite{W3C-DID2022}: identifier không phụ thuộc CA tập trung
% Verifiable Credentials (VC): chứng chỉ số thế hệ mới (bằng lái, văn bằng,...)
% Blockchain-based anchoring: DID document lưu trên distributed ledger
% Vẫn còn nhiều thách thức: UX, governance, key recovery khi mất private key
